\documentclass[11pt,a4paper]{report}

% --- Geometría de página ---
\usepackage[margin=2.5cm]{geometry}

% --- Espaciado de títulos (report deja demasiado aire antes/después) ---
\usepackage{titlesec}
\titleformat{\chapter}[hang]{\normalfont\Huge\bfseries}{\thechapter}{1em}{}
\titlespacing*{\chapter}{0pt}{0pt}{20pt}

% --- Idioma y codificación ---
\usepackage[utf8]{inputenc}
\usepackage[T1]{fontenc}
\usepackage[spanish,es-noindentfirst,es-noshorthands]{babel}
\usepackage{lmodern}

% --- Matemáticas y unidades ---
\usepackage{amsmath}
\usepackage{amssymb}
\usepackage{siunitx}
\sisetup{output-decimal-marker = {,}, group-separator = {.}}

% --- Figuras ---
\usepackage{tikz}
\usetikzlibrary{arrows.meta, positioning, calc, decorations.pathmorphing, patterns}
\usepackage{graphicx}
\usepackage{caption}

% --- Tablas ---
\usepackage{booktabs}
\usepackage{array}

% --- Código y monoespaciado ---
\usepackage{listings}
\usepackage{xcolor}
\lstset{
  basicstyle=\ttfamily\small,
  breaklines=true,
  frame=single,
  backgroundcolor=\color{gray!8},
  columns=fullflexible,
  keepspaces=true
}

% --- Cajas de nota didáctica ---
\usepackage[most]{tcolorbox}
\newtcolorbox{nota}[1][]{
  colback=blue!5!white, colframe=blue!40!black,
  fonttitle=\bfseries, title={Nota}, #1
}
\newtcolorbox{ejemplo}[1][]{
  colback=green!5!white, colframe=green!30!black,
  fonttitle=\bfseries, title={Ejemplo resuelto}, #1
}
\newtcolorbox{trampa}[1][]{
  colback=orange!5!white, colframe=orange!60!black,
  fonttitle=\bfseries, title={Trampa del hardware}, #1
}

% --- Referencias cruzadas y bibliografía ---
\usepackage{hyperref}
\hypersetup{
  colorlinks=true,
  linkcolor=blue!50!black,
  citecolor=green!40!black,
  urlcolor=blue!50!black,
  pdftitle={Fundamentos de telemetría LoRa},
  pdfauthor={Proyecto lorasync}
}

\title{Fundamentos de telemetría LoRa\\[0.3em]
  \large De la onda de radio al dato en el CSV}
\author{Proyecto \texttt{lorasync}}
\date{\today}

\begin{document}

\maketitle

\begin{abstract}
\noindent
Este documento explica, desde cero y con un tono didáctico, los fundamentos técnicos que
sostienen el sistema de telemetría LoRa del proyecto \texttt{lorasync}: por qué se eligieron
sus parámetros de radio, cómo funciona la modulación LoRa, por qué el espectro está regulado
y qué banda es legal en Colombia. Cada concepto se presenta en tres pasos: \emph{qué es}
(desde cero), \emph{por qué importa aquí} (aplicado a este proyecto) y \emph{dónde aparece en
el código}. Las partes restantes, que dependen del código ya escrito, se añaden en fases
posteriores del proyecto.
\end{abstract}

\tableofcontents

\part{Fundamentos de radio}
\label{part:radio}

\chapter{Qué es una onda de radio}

\section{Qué es}

Una onda de radio es una oscilación del campo electromagnético que viaja por el espacio a la
velocidad de la luz, $c \approx \SI{3e8}{m/s}$. Igual que una ola en el agua, una onda de radio
tiene una \emph{longitud de onda} $\lambda$ (la distancia entre dos crestas) y una
\emph{frecuencia} $f$ (cuántas crestas pasan por un punto fijo cada segundo). Las dos están
ligadas por una relación muy simple:

\begin{equation}
  c = f \cdot \lambda
  \label{eq:onda}
\end{equation}

\begin{figure}[h]
  \centering
  \begin{tikzpicture}[scale=1.0]
    \draw[-{Latex}] (0,0) -- (8.5,0) node[right] {tiempo};
    \draw[-{Latex}] (0,-1.3) -- (0,1.3) node[above] {amplitud};
    \draw[thick, blue!60!black, domain=0:8, samples=200]
      plot (\x, {sin(\x*180)});
    % marca de un periodo
    \draw[<->, thick] (0,-1.1) -- (2,-1.1) node[midway, below] {$T = 1/f$};
    \draw[dashed, gray] (0,1) -- (0,-1.3);
    \draw[dashed, gray] (2,1) -- (2,-1.3);
  \end{tikzpicture}
  \caption{Una onda senoidal: el periodo $T$ es el tiempo entre dos crestas, y la frecuencia
    es $f = 1/T$.}
  \label{fig:onda-seno}
\end{figure}

Cuando decimos que un módulo transmite a \textbf{\SI{920}{MHz}}, decimos que el campo
electromagnético oscila \num{920000000} veces por segundo. Con la ecuación~\eqref{eq:onda}, la
longitud de onda correspondiente es

\[
  \lambda = \frac{c}{f} = \frac{\SI{3e8}{m/s}}{\SI{920e6}{Hz}} \approx \SI{0.326}{m},
\]

es decir, unos 33 cm entre cresta y cresta. Este número no es anecdótico: el tamaño físico de
una antena eficiente está ligado a una fracción de $\lambda$ (típicamente $\lambda/4$ o
$\lambda/2$), razón por la cual las antenas para esta banda miden unos pocos centímetros.

\begin{nota}
La radio no ``sabe'' qué frecuencia debería usar cada aplicación. La asignación de bandas
(\SI{920}{MHz} para este proyecto, \SI{2.4}{GHz} para wifi, \SI{100}{MHz} para FM) es una
decisión humana, administrativa, que se estudia en la Parte~\ref{part:regulacion}.
\end{nota}

\section{Por qué importa aquí}

Todos los cálculos del sistema \texttt{lorasync} —presupuesto de enlace, tamaño de antena,
adaptación de impedancia— parten de un único número: la frecuencia de trabajo. El proyecto usa
el canal 70, que corresponde a \SI{920}{MHz} (ver la fórmula de canalización en la
Parte~\ref{part:regulacion}). Todo lo que sigue en este capítulo —decibelios, potencia,
ganancia, presupuesto de enlace— se calcula \emph{para esa frecuencia concreta}.

\section{Dónde aparece en el código}

La frecuencia no se transmite como tal; se configura indirectamente mediante el canal
(\texttt{AT+TXCH=70}, \texttt{AT+RXCH=70}), que el módulo traduce internamente a
\SI{920}{MHz}. La validación de que el canal esté en el rango legal vive en
\texttt{common/config.py} (Fase~1).

\chapter{Decibelios y dBm}

\section{Qué es}

Las potencias de radio abarcan un rango enorme: un transmisor puede emitir \SI{100}{mW}
mientras que un receptor detecta señales de una billonésima de vatio. Trabajar con esos
números en vatios es incómodo, así que la ingeniería de radio usa una escala
\textbf{logarítmica}: el decibelio (dB).

\begin{equation}
  \text{ganancia en dB} = 10 \cdot \log_{10}\!\left(\frac{P_2}{P_1}\right)
  \label{eq:db}
\end{equation}

La ventaja de una escala logarítmica es que convierte multiplicaciones en sumas. Doblar una
potencia (multiplicarla por 2) equivale siempre a sumar $10 \cdot \log_{10}(2) \approx
\SI{3}{dB}$, sin importar el punto de partida:

\begin{ejemplo}
Si un amplificador dobla la potencia de una señal, su ganancia es
$10\log_{10}(2) = \SI{3.01}{dB} \approx \SI{3}{dB}$.
Si la potencia se multiplica por 10, la ganancia es $10\log_{10}(10) = \SI{10}{dB}$.
Si se divide entre 2 (la mitad), la ``ganancia'' es $\SI{-3}{dB}$: una pérdida de \SI{3}{dB}.
\end{ejemplo}

El \textbf{dBm} es un caso particular: un decibelio \emph{referenciado a \SI{1}{mW}}.

\begin{equation}
  P_{\text{dBm}} = 10 \cdot \log_{10}\!\left(\frac{P_{\text{mW}}}{1\,\text{mW}}\right)
  \label{eq:dbm}
\end{equation}

\begin{table}[h]
  \centering
  \begin{tabular}{@{}rl@{}}
    \toprule
    Potencia & Equivalente \\
    \midrule
    \SI{0}{dBm}  & \SI{1}{mW} \\
    \SI{10}{dBm} & \SI{10}{mW} \\
    \SI{20}{dBm} & \SI{100}{mW} \\
    \SI{-118}{dBm} & $\approx \SI{1.6e-15}{W}$ (sensibilidad típica del receptor) \\
    \bottomrule
  \end{tabular}
  \caption{La escala dBm hace manejables números que en vatios serían absurdamente pequeños.}
\end{table}

\begin{nota}
Por ser una escala logarítmica \emph{referenciada}, los dBm \textbf{se suman y se restan}, no
se multiplican. Sumar dos etapas de ganancia en dB equivale a multiplicar las potencias reales.
Esto es exactamente lo que se aprovecha en el presupuesto de enlace (sección~\ref{sec:presupuesto}).
\end{nota}

\section{Por qué importa aquí}

Todo el presupuesto de enlace del proyecto —potencia de salida, ganancia de antena, pérdidas
por trayecto, sensibilidad del receptor— se expresa en dB y dBm precisamente porque así las
cuentas se reducen a sumas y restas en vez de productos y cocientes de números con muchos
ceros.

\section{Dónde aparece en el código}

El comando \texttt{AT+PWR=10} configura la potencia de transmisión en dBm directamente
(\num{10} a \num{22} dBm según el firmware). El RSSI que reporta el módulo
(\texttt{AT+RSSI=1}) también se expresa en dBm, y es el valor que se separa del final de cada
trama en \texttt{common/radio.py} (Fase~1) y se guarda en la columna \texttt{rssi} del CSV
(Fase~3).

\chapter{Potencia, ganancia de antena y PIRE}

\section{Qué es}

Tres cantidades determinan cuánta energía efectiva sale de un transmisor hacia el aire:

\begin{itemize}
  \item \textbf{Potencia de transmisión} ($P_{TX}$): la energía que el módulo entrega a la
    antena, en dBm. En este proyecto, \SI{10}{dBm}.
  \item \textbf{Ganancia de antena} ($G$): cuánto \emph{concentra} la antena esa energía en
    ciertas direcciones en vez de irradiarla por igual en todas. Se mide en dBi (decibelios
    respecto a una antena isotrópica ideal, que irradia igual en todas direcciones). Una antena
    típica de \SI{915}{MHz} tiene \SI{2}{dBi} a \SI{3}{dBi} de ganancia.
  \item \textbf{PIRE} (Potencia Isotrópica Radiada Equivalente, EIRP en inglés): la potencia
    que \emph{parecería} tener un transmisor isotrópico ideal para producir la misma intensidad
    de señal en la dirección de máxima ganancia de la antena real.
\end{itemize}

\begin{equation}
  \text{PIRE}_{\text{dBm}} = P_{TX,\text{dBm}} + G_{\text{dBi}}
  \label{eq:pire}
\end{equation}

Gracias a la escala logarítmica, calcular la PIRE es una simple suma.

\section{Por qué importa aquí}

Con $P_{TX} = \SI{10}{dBm}$ y una antena de $G = \SI{2}{dBi}$ en el nodo transmisor:

\[
  \text{PIRE} = 10\,\text{dBm} + 2\,\text{dBi} = 12\,\text{dBm}.
\]

Este valor es el punto de partida del presupuesto de enlace de la siguiente sección. La
ganancia de la antena \emph{receptora} también entra en la cuenta, porque una antena
receptora concentra igual de bien la energía que le llega.

\section{Dónde aparece en el código}

Ninguna de estas cantidades se calcula en tiempo de ejecución: son parámetros de diseño fijos
(potencia por \texttt{AT+PWR}, ganancia de antena elegida al comprar el hardware) que
justifican por qué \SI{10}{dBm} —y no la potencia máxima de \SI{22}{dBm}— es suficiente para
este proyecto, como se ve a continuación.

\chapter{Presupuesto de enlace}
\label{sec:presupuesto}

\section{Qué es}

Un \emph{presupuesto de enlace} (\emph{link budget}) es la contabilidad completa, en dB, de
toda la energía que gana o pierde una señal entre el transmisor y el receptor. Se parte de la
PIRE, se resta todo lo que la señal pierde en el camino, y se compara el resultado con la
sensibilidad del receptor: el nivel mínimo de señal que el receptor todavía puede decodificar.
Si lo que llega supera la sensibilidad, el enlace funciona; la diferencia es el
\textbf{margen de enlace}, que actúa de colchón frente a imprevistos (lluvia, obstáculos,
desvanecimientos).

La mayor pérdida en un enlace de línea de vista al aire libre es la \textbf{pérdida de espacio
libre} (FSPL, \emph{Free-Space Path Loss}): la dilución de la energía porque se reparte sobre
una esfera cada vez mayor a medida que la onda se aleja de la antena. La fórmula estándar,
para $d$ en kilómetros y $f$ en megahercios, es:

\begin{equation}
  \text{FSPL}_{\text{dB}} = 32.45 + 20 \log_{10}(f_{\text{MHz}}) + 20 \log_{10}(d_{\text{km}})
  \label{eq:fspl}
\end{equation}

El $20\log_{10}$ en ambos términos refleja que la pérdida crece con el \emph{cuadrado} de la
frecuencia y con el \emph{cuadrado} de la distancia: doblar la distancia cuesta
$20\log_{10}(2) \approx \SI{6}{dB}$ adicionales de pérdida.

\section{Por qué importa aquí: los 600 m resueltos}

El proyecto exige cobertura fiable a \SI{600}{m} en el canal 70 (\SI{920}{MHz}), con
$P_{TX} = \SI{10}{dBm}$, antenas de \SI{2}{dBi} en ambos extremos y una sensibilidad de
receptor de $\approx \SI{-118}{dBm}$ para SF7 a \SI{500}{kHz} de ancho de banda (el porqué de
esta cifra se explica en la Parte~\ref{part:lora}).

\begin{ejemplo}
\textbf{Paso 1 — Pérdida de espacio libre a 600 m y 920 MHz:}
\begin{align*}
  \text{FSPL} &= 32.45 + 20\log_{10}(920) + 20\log_{10}(0.6) \\
              &= 32.45 + 59.28 - 4.44 \\
              &= \SI{87.3}{dB}
\end{align*}

\textbf{Paso 2 — Potencia recibida:}
\begin{align*}
  P_{RX} &= P_{TX} + G_{TX} + G_{RX} - \text{FSPL} \\
         &= 10\,\text{dBm} + 2\,\text{dBi} + 2\,\text{dBi} - 87.3\,\text{dB} \\
         &= \SI{-73.3}{dBm}
\end{align*}

\textbf{Paso 3 — Margen de enlace:}
\begin{align*}
  \text{Margen} &= P_{RX} - \text{Sensibilidad} \\
                &= -73.3\,\text{dBm} - (-118\,\text{dBm}) \\
                &\approx \SI{45}{dB}
\end{align*}
\end{ejemplo}

Un margen de \SI{45}{dB} es muy holgado para \SI{600}{m}: en la práctica, este enlace
funcionaría a distancias mucho mayores en condiciones de línea de vista despejada. El margen
existe precisamente para absorber lo que la fórmula ideal no modela: vegetación, paredes,
desvanecimiento por multitrayecto y el hecho de que las antenas casi nunca están perfectamente
alineadas.

\begin{figure}[h]
  \centering
  \begin{tikzpicture}[node distance=1.3cm, font=\small, every node/.style={align=center}]
    \node[draw, rounded corners, fill=blue!8] (tx) {PIRE\\ \SI{12}{dBm}};
    \node[draw, rounded corners, fill=red!8, right=2.6cm of tx] (fspl) {$-$ FSPL\\ \SI{87.3}{dB}};
    \node[draw, rounded corners, fill=blue!8, right=2.6cm of fspl] (rx) {$P_{RX}$\\ \SI{-73.3}{dBm}};
    \draw[-{Latex}, thick] (tx) -- (fspl);
    \draw[-{Latex}, thick] (fspl) -- (rx);
    \node[below=0.6cm of rx, draw, rounded corners, fill=green!10] (sens) {Sensibilidad\\ \SI{-118}{dBm}};
    \draw[<->, thick, green!30!black] ($(rx.south) + (0,-0.05)$) -- (sens.north)
      node[midway, right] {margen $\approx \SI{45}{dB}$};
  \end{tikzpicture}
  \caption{Presupuesto de enlace a \SI{600}{m}: de la PIRE a la potencia recibida, y el margen
    frente a la sensibilidad del receptor.}
\end{figure}

\section{Adaptación de antena}

Una antena está diseñada para una frecuencia central, pero sigue funcionando razonablemente
bien en frecuencias cercanas: su comportamiento (impedancia, patrón de radiación) varía de
forma suave, no abrupta. Una antena diseñada para \SI{868}{MHz}, usada a \SI{920}{MHz}, tiene
un desajuste de apenas

\[
  \frac{920 - 868}{868} \approx 0.06 \approx \SI{6}{\percent},
\]

que en la práctica se traduce en una pérdida de adaptación de solo \textasciitilde\SI{0.2}{dB}
a \SI{0.5}{dB} —imperceptible frente a un margen de enlace de \SI{45}{dB}—. Por eso las
antenas de \SI{868}{MHz} que ya existían sirven para arrancar el proyecto, aunque a futuro
conviene reemplazarlas por antenas centradas en \SI{915}{MHz} para maximizar el rendimiento y
evitar cualquier duda regulatoria sobre el equipo irradiante.

\section{Dónde aparece en el código}

\texttt{common/airtime.py} (Fase~1) no calcula el presupuesto de enlace —es una decisión de
diseño ya tomada, no un valor que el software deba verificar en tiempo real—, pero sí depende
de los mismos parámetros de radio (SF, BW) que determinan la sensibilidad usada aquí. La
herramienta \texttt{tools/linkcheck.py} (Fase~5) mide el RSSI real en campo para contrastarlo
con el $P_{RX} \approx \SI{-73.3}{dBm}$ calculado en esta sección.

\part{Cómo funciona LoRa}
\label{part:lora}

\chapter{Qué es modular una señal}

\section{Qué es}

Una antena por sí sola solo puede emitir una onda a una frecuencia fija; para transportar
información hay que \textbf{modular}: variar alguna propiedad de esa onda (amplitud,
frecuencia o fase) de forma que el receptor pueda reconstruir los bits originales a partir de
esas variaciones. AM varía la amplitud, FM varía la frecuencia; los esquemas digitales modernos
—incluido LoRa— varían combinaciones de frecuencia y fase para codificar símbolos, no solo bits
sueltos.

\section{Por qué importa aquí}

LoRa no usa un esquema de modulación clásico como FSK o PSK puros: usa \textbf{espectro
ensanchado por chirp} (\emph{Chirp Spread Spectrum}, CSS), diseñado específicamente para
maximizar el alcance a costa de velocidad, que es exactamente el compromiso que necesita un
sensor remoto que envía pocos bytes por segundo pero debe llegar lejos con un transmisor de
baja potencia.

\chapter{Espectro ensanchado por chirp}

\section{Qué es}

Un \emph{chirp} es un tono cuya frecuencia sube (o baja) linealmente en el tiempo, barriendo
todo el ancho de banda del canal durante la duración de un símbolo. Es el mismo fenómeno que un
silbido que empieza grave y termina agudo.

\begin{nota}[title={Analogía}]
Imagina que en vez de gritar una sola nota fija (como haría una modulación de amplitud
clásica), silbas deslizando el tono de grave a agudo. Ese deslizamiento completo, con una
velocidad de barrido determinada, \emph{es} un símbolo. LoRa codifica información desplazando
el \emph{punto de partida} de ese barrido: un silbido que empieza a medio camino y da la vuelta
codifica un símbolo distinto de uno que empieza al principio.
\end{nota}

\begin{figure}[h]
  \centering
  \begin{tikzpicture}[scale=1.0, font=\footnotesize]
    \draw[-{Latex}] (0,0) -- (6.5,0) node[right] {tiempo};
    \draw[-{Latex}] (0,0) -- (0,3.2) node[above] {frecuencia};
    % up-chirp completo (símbolo 0)
    \draw[thick, blue!60!black] (0,0.3) -- (2,3);
    \draw[thick, blue!60!black, dashed] (2,3) -- (2,0);
    \draw[thick, blue!60!black] (2,0) -- (4,2.7);
    \node[blue!60!black, below, text width=2.6cm, align=center] at (1.6,-0.15)
      {símbolo 0\\(arranca en 0)};
    % chirp con offset (símbolo desplazado)
    \draw[thick, red!60!black] (4,1.6) -- (4.9,3);
    \draw[thick, red!60!black, dashed] (4.9,3) -- (4.9,0);
    \draw[thick, red!60!black] (4.9,0) -- (6,1.6);
    \node[red!60!black, below, text width=2.4cm, align=center] at (5.3,-0.15)
      {símbolo\\con offset};
  \end{tikzpicture}
  \caption{Un chirp barre linealmente toda la banda; el punto donde empieza el barrido
    (y da la vuelta) codifica el valor del símbolo. Cada línea discontinua marca un
    ``salto'' de vuelta a la frecuencia mínima al llegar al máximo del canal.}
\end{figure}

Por qué esto ayuda al alcance: el receptor no necesita detectar un chirp aislado con precisión
absoluta de amplitud (como sí exige FSK/PSK clásicos); necesita \emph{correlacionar} la señal
recibida con la forma conocida del chirp. Esa correlación acumula energía de todo el símbolo, lo
que permite decodificar señales que están \emph{por debajo del nivel de ruido} —algo imposible
con modulaciones de banda estrecha convencionales—.

\section{Factor de dispersión: el compromiso alcance/velocidad}

\subsection*{Qué es}

El \textbf{factor de dispersión} (\emph{Spreading Factor}, SF) determina cuántos \emph{bits por
símbolo} se transmiten y, por tanto, cuánto dura cada chirp: $2^{SF}$ posibles puntos de
arranque distintos, es decir, $SF$ bits de información por símbolo. Un SF más alto significa
más puntos de arranque posibles (más información por símbolo) pero también un barrido más
lento, así que cada símbolo tarda más tiempo en transmitirse.

\begin{equation}
  T_{\text{sym}} = \frac{2^{SF}}{BW}
  \label{eq:tsym}
\end{equation}

\subsection*{Por qué importa aquí}

Un SF más alto extiende la energía de cada bit sobre un chirp más largo, lo que mejora la
\textbf{sensibilidad} del receptor (puede detectar señales más débiles) a costa de una
\textbf{tasa de datos} menor. Este proyecto usa \textbf{SF7}, el más bajo (y más rápido) del
rango típico (7 a 12), porque el presupuesto de enlace de la sección~\ref{sec:presupuesto}
mostró un margen de \SI{45}{dB} incluso con SF7: no hace falta sacrificar velocidad para ganar
alcance que ya sobra.

\begin{table}[h]
  \centering
  \small
  \begin{tabular}{@{}lrrl@{}}
    \toprule
    SF & Símbolos/s @ 500 kHz & Sensibilidad aprox. & Compromiso \\
    \midrule
    7  & \num{3906}  & $\approx \SI{-118}{dBm}$ & rápido, alcance moderado \\
    9  & \num{977}   & $\approx \SI{-124}{dBm}$ & más lento, más alcance \\
    12 & \num{122}   & $\approx \SI{-131}{dBm}$ & muy lento, máximo alcance \\
    \bottomrule
  \end{tabular}
  \caption{A mayor SF, mayor sensibilidad (más alcance) pero menor tasa de símbolos
    (más tiempo de aire por trama). Valores de sensibilidad aproximados para BW = 500 kHz.}
\end{table}

\section{Ancho de banda y corrección de errores}

\subsection*{Qué es}

El \textbf{ancho de banda} (BW) es el rango de frecuencias que ocupa cada chirp: cuanto más
ancho, más rápido se puede barrer (menos tiempo por símbolo, ecuación~\eqref{eq:tsym}), pero
también entra más ruido de banda ancha al receptor, lo que reduce ligeramente la sensibilidad.
Este proyecto usa \textbf{\SI{500}{kHz}}, el máximo disponible en el módulo, para minimizar el
tiempo de aire.

La \textbf{tasa de codificación} (\emph{Coding Rate}, CR) añade bits redundantes de corrección
de errores hacia adelante (FEC): con CR = 4/5, por cada 4 bits de datos se transmite 1 bit de
paridad. Esto permite al receptor corregir errores de bit sin pedir retransmisión, a costa de
transmitir algo más de información por trama útil. El proyecto usa \textbf{4/5}, la tasa mínima
de sobrecarga (frente a 4/6, 4/7 o 4/8, más robustas pero más lentas).

\subsection*{Por qué importa aquí}

BW alto y CR mínimo son ambos decisiones de \emph{minimizar el tiempo de aire}, coherentes con
la elección de SF7: el margen de enlace de \SI{45}{dB} permite ser agresivo en velocidad sin
poner en riesgo la fiabilidad del enlace a \SI{600}{m}.

\chapter{El tiempo de aire, término a término}
\label{sec:airtime}

\section{Qué es}

El \textbf{tiempo de aire} (\emph{time on air}) es cuánto dura, en el aire, la transmisión
completa de una trama: preámbulo más carga útil. Es el número que determina cuántas tramas por
segundo caben en el canal y, por extensión, cuántos nodos puede soportar el sistema —el tema de
la Parte~\ref{part:regulacion} y del diseño TDMA de la Fase~4—.

La fórmula estándar de LoRa (documentada por Semtech) tiene dos partes: el preámbulo y la carga
útil, cada una expresada en número de símbolos y luego convertida a tiempo con $T_{\text{sym}}$
de la ecuación~\eqref{eq:tsym}.

\subsection*{Preámbulo}

\begin{equation}
  T_{\text{preámbulo}} = (n_{\text{pre}} + 4.25) \cdot T_{\text{sym}}
  \label{eq:tpre}
\end{equation}

$n_{\text{pre}}$ es el número de símbolos de preámbulo configurados (típicamente 8: una
secuencia conocida que el receptor usa para sincronizarse). El término fijo $4.25$ cubre la
palabra de sincronización y el arranque de la decodificación, y es parte de la especificación
de Semtech, no un parámetro configurable.

\subsection*{Carga útil}

\begin{equation}
  n_{\text{payload}} = 8 + \max\!\left(
    \left\lceil \frac{8 \cdot PL - 4 \cdot SF + 28 + 16}{4 \cdot SF} \right\rceil
    \cdot (CR + 4),\ 0 \right)
  \label{eq:npayload}
\end{equation}

\begin{equation}
  T_{\text{trama}} = T_{\text{preámbulo}} + n_{\text{payload}} \cdot T_{\text{sym}}
  \label{eq:ttrama}
\end{equation}

donde $PL$ es el tamaño de la carga útil en bytes, y $CR$ es la tasa de codificación expresada
como el número de bits de paridad por cada 4 bits de datos (CR = 1 para 4/5, CR = 4 para 4/8).
El término $+16$ dentro del numerador corresponde al CRC de 16 bits que valida la integridad de
la trama a nivel de radio (independiente del CRC16 propio de la capa de aplicación, que se
implementa en la Fase~2); el término $+28$ y el $8$ inicial son constantes de la
especificación LoRa relacionadas con el encabezado explícito.

\section{Por qué importa aquí: los 76,9 ms resueltos}

El proyecto agrupa 10 muestras por trama. Con el formato de trama de la Fase~2 (cabecera +
10 muestras de \SI{18}{B} + CRC16, ver Supuestos declarados en el plan), el tamaño resultante
es $PL = \SI{194}{B}$. Con SF7, BW = \SI{500}{kHz} y CR = 1 (4/5):

\begin{ejemplo}
\textbf{Paso 1 — Duración de un símbolo:}
\[
  T_{\text{sym}} = \frac{2^7}{500\,000} = \SI{0.256}{ms}
\]

\textbf{Paso 2 — Preámbulo (8 símbolos configurados):}
\[
  T_{\text{preámbulo}} = (8 + 4.25) \cdot 0.256\,\text{ms} = \SI{3.136}{ms}
\]

\textbf{Paso 3 — Símbolos de carga útil, con $PL = 194$, $SF = 7$, $CR = 1$:}
\begin{align*}
  \frac{8 \cdot 194 - 4 \cdot 7 + 28 + 16}{4 \cdot 7}
    &= \frac{1552 - 28 + 28 + 16}{28} = \frac{1568}{28} = 56 \\
  n_{\text{payload}} &= 8 + \lceil 56 \rceil \cdot (1 + 4) = 8 + 56 \cdot 5 = 288
\end{align*}

\textbf{Paso 4 — Tiempo de trama:}
\begin{align*}
  T_{\text{trama}} &= 3.136\,\text{ms} + 288 \cdot 0.256\,\text{ms} \\
    &= 3.136 + 73.728 = \SI{76.864}{ms} \approx \SI{76.9}{ms}
\end{align*}
\end{ejemplo}

Este es exactamente el valor que reproduce \texttt{common/airtime.py} en la Fase~1, y el que
sostiene el diseño del slot TDMA de \SI{136.9}{ms} (76,9 ms de trama + 60 ms de guarda) en la
Fase~4.

\section{Por qué un SF alto no siempre es mejor}

Es tentador pensar que, si SF7 ya funciona con \SI{45}{dB} de margen, usar SF12 —el más
robusto— sería ``más seguro''. En este proyecto sería contraproducente: la
ecuación~\eqref{eq:tsym} muestra que $T_{\text{sym}}$ crece \emph{exponencialmente} con SF, no
linealmente. Pasar de SF7 a SF12 multiplica $T_{\text{sym}}$ por $2^5 = 32$, lo que dispararía
el tiempo de aire de cada trama muy por encima de lo que 30 nodos pueden compartir en un canal
sin colisionar constantemente (ver el diseño de superframe de la Fase~4). El SF correcto no es
el más robusto en abstracto, sino el más bajo que todavía deja margen de enlace suficiente para
la distancia real del despliegue —exactamente el criterio que usó la sección~\ref{sec:presupuesto}
para justificar SF7—.

\section{Dónde aparece en el código}

\texttt{common/airtime.py} (Fase~1) implementa las ecuaciones~\eqref{eq:tsym},
\eqref{eq:tpre}, \eqref{eq:npayload} y \eqref{eq:ttrama} literalmente, y su prueba
(\texttt{tests/test\_airtime.py}) verifica el resultado de \SI{76.9}{ms} calculado aquí. El
mismo módulo calcula la duración de slot y la capacidad del canal que sostienen el diseño
TDMA de la Fase~4.

\part{El espectro como recurso legal}
\label{part:regulacion}

\chapter{Por qué se regula el espectro}

\section{Qué es}

El espectro radioeléctrico es un recurso \textbf{compartido y finito}: dos transmisores en la
misma frecuencia, al mismo tiempo y en el mismo lugar, se interfieren mutuamente y ninguno se
puede decodificar. A diferencia de un cable, donde añadir más usuarios es cuestión de tender más
cobre, el espectro no se puede ``fabricar'': todo el que existe ya está ahí, y su uso debe
coordinarse por ley para que servicios críticos —aviación, emergencias, telefonía celular
licenciada— no sufran interferencia de dispositivos no autorizados.

Por eso cada país tiene una autoridad reguladora que decide qué franjas de frecuencia se
reservan para qué uso, y bajo qué condiciones técnicas (potencia máxima, ancho de banda, ciclo
de trabajo) se puede transmitir en cada una.

\section{Qué son las bandas ISM}

Las bandas \textbf{ISM} (\emph{Industrial, Scientific and Medical}) son franjas de espectro
que la mayoría de países reserva para uso \textbf{libre y sin licencia individual}, siempre que
el equipo cumpla ciertas condiciones técnicas (potencia máxima, tipo de modulación). Es el
régimen bajo el cual operan wifi (\SI{2.4}{GHz}), Bluetooth, y también las bandas sub-GHz que
usa LoRa. ``Libre'' no significa ``sin reglas'': significa que no hace falta pedir una licencia
de operación individual, pero el equipo sí debe respetar los límites técnicos de la banda.

\section{Por qué importa aquí}

El proyecto \texttt{lorasync} depende enteramente de que su banda de operación sea, en efecto,
de uso libre en Colombia. Transmitir en una banda licenciada sin autorización —aunque sea con
poca potencia— es una infracción, no un simple detalle técnico.

\chapter{Qué es la ANE}

La \textbf{Agencia Nacional del Espectro} (ANE) es la entidad colombiana que administra,
vigila y controla el uso del espectro radioeléctrico en el país. Publica el \textbf{Cuadro
Nacional de Atribución de Bandas de Frecuencias} (CNABF), el documento que determina, banda por
banda, qué servicios están autorizados y bajo qué condiciones. Es la referencia oficial usada
para las cifras de esta sección~\cite{ane_cnabf}.

\chapter{Por qué 915--928 MHz sí y 868 MHz no}

\section{Qué es}

En Europa, la banda sub-GHz de uso libre para IoT es \SI{863}{}--\SI{870}{MHz}, con
\SI{868}{MHz} como frecuencia de referencia habitual —de ahí que muchísima documentación y
hardware genérico de LoRa (incluidas antenas) esté etiquetado ``868 MHz''—. Colombia
\textbf{no} sigue la atribución europea en esa franja: la ANE clasifica
\SI{851}{}--\SI{915}{MHz} como espectro celular licenciado (usado por los operadores móviles),
y reserva la banda de uso libre en \SI{915}{}--\SI{928}{MHz}, alineada con el modelo
norteamericano de banda ISM de 900 MHz~\cite{ttn_au915}.

\begin{figure}[h]
  \centering
  \begin{tikzpicture}[scale=1.0, font=\footnotesize, every node/.style={align=center}]
    % eje
    \draw[-{Latex}] (0,0) -- (10.5,0) node[right, font=\footnotesize] {frecuencia (MHz)};
    % banda restringida (celular licenciado)
    \fill[red!25] (0.5,0.3) rectangle (7,1.7);
    \draw (0.5,0.3) rectangle (7,1.7);
    \node[text width=5.6cm] at (3.75,1.0) {celular licenciado\\(851--915 MHz)\\\textbf{prohibido}};
    % banda libre
    \fill[green!25] (7,0.3) rectangle (9,1.7);
    \draw (7,0.3) rectangle (9,1.7);
    \node at (8,1.0) {\textbf{915--928}\\libre};
    % ticks y etiquetas del eje
    \foreach \x/\lbl in {0.5/851, 7/915, 9/928} {
      \draw (\x,0.25) -- (\x,0.35);
      \node[below, font=\scriptsize] at (\x,0.15) {\lbl};
    }
    % marca 868 (bajo el eje, separada de las etiquetas de los ticks)
    \draw[dashed, blue!60!black, thick] (4.55,-0.9) -- (4.55,1.9);
    \node[blue!60!black, font=\scriptsize, text width=3.2cm] at (4.55,-1.15)
      {868 MHz (Europa)\\\textbf{no usar en CO}};
    % marca canal 70 -> 920 MHz (sobre el eje)
    \draw[dashed, orange!70!black, thick] (7.9,-0.9) -- (7.9,2.1);
    \node[orange!70!black, font=\scriptsize, text width=3.2cm] at (7.9,2.35)
      {canal 70 = 920 MHz\\(usado aquí)};
  \end{tikzpicture}
  \caption{La banda de \SI{868}{MHz}, libre en Europa, cae dentro del espectro celular
    licenciado en Colombia. La banda libre colombiana equivalente es 915--928\,MHz.}
  \label{fig:banda-colombia}
\end{figure}

\section{Por qué importa aquí}

Confundir la convención europea con la colombiana es el error regulatorio más probable en un
proyecto que usa hardware genérico de LoRa (documentación, foros y valores de fábrica casi
siempre asumen \SI{868}{MHz}). El valor de fábrica del módulo Waveshare (canal 18) corresponde
precisamente a \SI{868}{MHz}, y \textbf{no debe usarse en Colombia}.

\section{La fórmula de canalización}

El firmware del módulo no acepta una frecuencia directamente: acepta un número de
\textbf{canal}, que se traduce a frecuencia con:

\begin{equation}
  f_{\text{MHz}} = 850 + \text{canal}
  \label{eq:canal}
\end{equation}

Los canales válidos para la banda libre colombiana son, por tanto, los que caen entre 915 y
928: $\text{canal} \in [65, 78]$. El proyecto usa el \textbf{canal 70}, que por la
ecuación~\eqref{eq:canal} da $850 + 70 = \SI{920}{MHz}$, cómodamente dentro del rango legal y
lejos de sus bordes.

\begin{trampa}
El valor de fábrica del módulo es el \textbf{canal 18}, que da $850 + 18 = \SI{868}{MHz}$: la
frecuencia europea, ilegal en Colombia. Cualquier módulo recién desempacado o reseteado a
valores de fábrica transmite fuera de banda hasta que se reconfigura explícitamente.
\end{trampa}

\chapter{Ciclo de trabajo y convivencia}

\section{Qué es}

El \textbf{ciclo de trabajo} (\emph{duty cycle}) es el porcentaje de tiempo que un transmisor
tiene permitido ocupar el canal. En Europa, la banda de \SI{868}{MHz} impone un límite estricto
de ciclo de trabajo (típicamente 1\,\% por sub-banda) como mecanismo regulatorio para forzar a
que muchos dispositivos puedan compartir el mismo canal sin coordinación explícita. El modelo
norteamericano (FCC Part 15.247), del que deriva la banda libre de 900 MHz en Colombia, no
impone un límite de ciclo de trabajo tan estricto, pero exige en su lugar \textbf{ancho de
banda mínimo} o \textbf{salto de frecuencia} (\emph{frequency hopping}), como mecanismo
alternativo para evitar que un solo dispositivo monopolice una franja estrecha del
espectro~\cite{fcc_15247}.

\section{Por qué importa aquí}

Los módulos Waveshare de este proyecto son de \textbf{canal fijo}: no implementan salto de
frecuencia. Esto significa que no pueden cumplir la variante ``frequency hopping'' del modelo
FCC. La variante que sí aplica es la de \textbf{ancho de banda mínimo mediante modulación
digital de banda ancha}: los \SI{500}{kHz} de ancho de banda elegidos para este proyecto no son
solo la opción de máxima capacidad (menor tiempo de aire, sección~\ref{sec:airtime}) —también
son la que mejor encaja con esta condición regulatoria—.

\begin{nota}
Esta interpretación es razonable para un prototipo de investigación con 2--30 nodos, pero
\textbf{no sustituye una verificación formal con la ANE} antes de cualquier despliegue
permanente o comercial. Es uno de los supuestos declarados explícitamente en el plan del
proyecto.
\end{nota}

Con solo 2--30 nodos transmitiendo tramas cortas de forma coordinada (ver el diseño TDMA de la
Fase~4), la ocupación real del canal —41,5\,\% en el escenario de 30 nodos— queda por debajo de
niveles que pondrían en riesgo la convivencia con otros usuarios de la misma banda libre.

\section{Dónde aparece en el código}

La validación del rango de canal (65--78) vive en \texttt{common/config.py} (Fase~1): el
sistema \textbf{se niega a arrancar} fuera de ese rango salvo que se pase explícitamente el
flag \texttt{allow\_out\_of\_band}, precisamente para que un error de configuración no termine
transmitiendo en \SI{868}{MHz} por accidente.


% Partes IV-VIII: se añaden en la Fase 5, cuando el código correspondiente
% (radio.py, frame.py, store.py, schedule.py, main.py) ya existe.
% \input{secciones/04-modulo}
% \input{secciones/05-canal}
% \input{secciones/06-integridad}
% \input{secciones/07-persistencia}
% \input{secciones/08-sistema}
% \input{secciones/09-apendices}

\bibliographystyle{plain}
\bibliography{referencias}

\end{document}
